% !TeX program = lualatex
\documentclass[12pt,aspectratio=169]{beamer}
\usepackage[utf8]{inputenc}
\usepackage[T1]{fontenc}
\usepackage{amsmath,amsfonts,amssymb}
\usepackage{graphicx}
\usepackage{tikz}
\usepackage{xcolor}
\usepackage{hyperref}
\usepackage{anyfontsize}
\usepackage{lmodern}

% ====== MAKE ALL FONTS SMALLER ======
\renewcommand{\tiny}{\fontsize{5}{6}\selectfont}
\renewcommand{\scriptsize}{\fontsize{6}{7}\selectfont}
\renewcommand{\footnotesize}{\fontsize{7}{8}\selectfont}
\renewcommand{\small}{\fontsize{8}{9}\selectfont}
\renewcommand{\normalsize}{\fontsize{9}{10}\selectfont}
\renewcommand{\large}{\fontsize{10}{11}\selectfont}
\renewcommand{\Large}{\fontsize{11}{13}\selectfont}
\renewcommand{\LARGE}{\fontsize{13}{15}\selectfont}
\renewcommand{\huge}{\fontsize{15}{17}\selectfont}
\renewcommand{\Huge}{\fontsize{18}{21}\selectfont}

% ====== COLOR THEME ======
\definecolor{redcorrect}{RGB}{200,30,30}
\definecolor{bluecorrect}{RGB}{30,80,200}
\definecolor{greencheck}{RGB}{0,150,50}
\definecolor{lightgray}{RGB}{245,245,245}
\definecolor{darkgray}{RGB}{50,50,50}
\definecolor{softyellow}{RGB}{255,248,200}
\definecolor{infoblue}{RGB}{230,240,252}
\definecolor{steppurple}{RGB}{150,50,200}
\definecolor{deepblue}{RGB}{20,60,150}
\definecolor{orangewarning}{RGB}{255,140,0}

\setbeamercolor{title}{fg=white,bg=redcorrect}
\setbeamercolor{frametitle}{fg=white,bg=redcorrect}
\setbeamercolor{block title}{fg=white,bg=bluecorrect}
\setbeamercolor{block body}{fg=darkgray,bg=lightgray}
\setbeamercolor{normal text}{fg=darkgray}
\usecolortheme{default}

% ====== CUSTOM COMMANDS ======
\newcommand{\correct}[1]{\textcolor{greencheck}{\textbf{\checkmark}} \textcolor{greencheck}{#1}}
\newcommand{\highlightred}[1]{\textcolor{redcorrect}{\textbf{#1}}}
\newcommand{\highlightblue}[1]{\textcolor{bluecorrect}{\textbf{#1}}}
\newcommand{\highlightorange}[1]{\textcolor{orangewarning}{\textbf{#1}}}
\newcommand{\tipbox}[1]{\fcolorbox{bluecorrect}{softyellow}{\parbox{0.88\textwidth}{\centering #1}}}
\newcommand{\steptitle}[1]{\textcolor{steppurple}{\textbf{#1}}}
\newcommand{\keypoint}[1]{\textcolor{deepblue}{\textbf{#1}}}
\newcommand{\warningbox}[1]{\fcolorbox{orangewarning}{yellow!15}{\parbox{0.88\textwidth}{\centering #1}}}
\newcommand{\redtitle}[1]{\textcolor{redcorrect}{\textbf{#1}}}

\newcommand{\stepbox}[1]{%
	\begin{center}
		\fcolorbox{darkgray}{white}{%
			\parbox{0.90\textwidth}{\vspace{0.1cm} #1 \vspace{0.1cm}}%
		}
	\end{center}
}

\begin{document}
	
	% ================================================================
	% SLIDE 1: TITLE
	% ================================================================
	\begin{frame}
		\centering
		\vspace{0.8cm}
		{\fontsize{20}{24}\selectfont \textbf{Money Laundering and Financial Crime}}\\[0.2cm]
		{\fontsize{15}{18}\selectfont Complete Tutorial Guide — All New Topics}\\[0.4cm]
		{\fontsize{12}{14}\selectfont Step-by-Step Learning for Every Subject}\\[0.4cm]
		{\fontsize{10}{12}\selectfont Compliance Training Department}
		\vspace{0.3cm}
		\rule{5cm}{1pt}
	\end{frame}
	
	% ================================================================
	% SLIDE 2: TUTORIAL 1 — DNFBP RISK MANAGEMENT
	% ================================================================
	\begin{frame}
		\frametitle{\fontsize{14}{17}\selectfont TUTORIAL 1: Managing DNFBP Risks}
		
		\small
		\begin{block}{\fontsize{11}{13}\selectfont Topic: DNFBP Risk Management Strategies}
			\footnotesize
			Designated Non-Financial Businesses and Professions (DNFBPs) present unique risks. Institutions need specific strategies to manage these risks effectively.
		\end{block}
		
		\vspace{0.03cm}
		
		\stepbox{\footnotesize
			\steptitle{Step 1:} Understand the nature of DNFBPs. These include lawyers, accountants, real estate agents, trust and company service providers, and dealers in precious metals/stones. They are vulnerable because they handle client money and can facilitate transactions. \\
			\steptitle{Step 2:} Learn the first strategy. \textbf{Creating a reputational risk or client selection committee} helps manage DNFBP risks. This committee evaluates whether to accept or retain high-risk clients. \\
			\steptitle{Step 3:} Learn the second strategy. \textbf{Using standardized sector-specific questionnaires, such as Wolfsberg-style documents} helps gather consistent information across different DNFBP sectors. \\
			\steptitle{Step 4:} Understand why these are correct. The committee provides \textbf{oversight and decision-making} for high-risk cases. The questionnaires provide \textbf{consistent data collection} across the organization. \\
			\steptitle{Step 5:} See what is NOT correct. \textbf{Defaulting to enhanced due diligence for all DNFBPs} is not efficient or effective. A risk-based approach is better.}
		
		\vspace{0.08cm}
		\tipbox{\footnotesize \textbf{CORE LESSON:} Manage DNFBP risks with a client selection committee and standardized questionnaires. Do not apply blanket enhanced due diligence to all DNFBPs.}
	\end{frame}
	
	% ================================================================
	% SLIDE 3: TUTORIAL 2 — FREE-TRADE ZONE RISKS
	% ================================================================
	\begin{frame}
		\frametitle{\fontsize{14}{17}\selectfont TUTORIAL 2: Free-Trade Zone Risks}
		
		\small
		\begin{block}{\fontsize{11}{13}\selectfont Topic: Money Laundering Risks in FTZs}
			\footnotesize
			Free-Trade Zones (FTZs) are high-risk environments for money laundering. The key risk factor is the lack of standardized inspection protocols.
		\end{block}
		
		\vspace{0.03cm}
		
		\stepbox{\footnotesize
			\steptitle{Step 1:} Understand what Free-Trade Zones are. FTZs are designated areas where goods can be imported, handled, and re-exported without customs intervention. They are designed to facilitate international trade. \\
			\steptitle{Step 2:} Recognize the key risk factor. The \textbf{absence of standardized inspection protocols for transshipped goods} most significantly increases money laundering risk in FTZs. \\
			\steptitle{Step 3:} Understand why this is the key risk. When goods pass through multiple jurisdictions without consistent inspection, criminals can: \\
			\quad - Misrepresent the nature of goods \\
			\quad - Overvalue or undervalue shipments \\
			\quad - Move money through trade-based laundering \\
			\steptitle{Step 4:} See how criminals exploit this. They can ship goods through FTZs to hide the true origin or destination of funds. The lack of inspection creates a gap. \\
			\steptitle{Step 5:} Apply the lesson. When assessing FTZ risk, focus on inspection protocols. Without consistent checks, the zone becomes a vulnerability.}
		
		\vspace{0.08cm}
		\tipbox{\footnotesize \textbf{CORE LESSON:} Free-Trade Zones are high-risk because goods move through multiple jurisdictions without standardized inspection. This creates opportunities for trade-based money laundering.}
	\end{frame}
	
	% ================================================================
	% SLIDE 4: TUTORIAL 3 — TCSP CLIENT ACCEPTANCE
	% ================================================================
	\begin{frame}
		\frametitle{\fontsize{14}{17}\selectfont TUTORIAL 3: TCSP Client Acceptance}
		
		\small
		\begin{block}{\fontsize{11}{13}\selectfont Topic: Trust and Company Service Providers}
			\footnotesize
			TCSPs must carefully evaluate clients. Multiple red flags indicate high risk and require decisive action.
		\end{block}
		
		\vspace{0.03cm}
		
		\stepbox{\footnotesize
			\steptitle{Step 1:} Understand the scenario. A TCSP is reviewing a client who wants to incorporate in a low-tax Caribbean jurisdiction. The client requests confidentiality for all shareholders. There is a connection to industries marked high-risk in the National Risk Assessment. \\
			\steptitle{Step 2:} Recognize the red flags. \textbf{Three red flags}: (1) Low-tax jurisdiction, (2) Complete shareholder confidentiality, (3) Connection to high-risk industries. This combination is highly suspicious. \\
			\steptitle{Step 3:} Understand the correct response. The TCSP should \textbf{decline immediately unless full beneficial owner information and risk justifications are disclosed and verified}. \\
			\steptitle{Step 4:} See what is NOT correct. \textbf{Referring to offshore local counsel and detaching from compliance responsibility} is wrong. The TCSP cannot delegate its compliance obligations. \\
			\steptitle{Step 5:} Apply the lesson. When red flags accumulate, demand full disclosure or decline the client. Do not pass the problem to someone else.}
		
		\vspace{0.08cm}
		\tipbox{\footnotesize \textbf{CORE LESSON:} When a client has multiple red flags — high-risk jurisdiction, confidentiality requests, and connection to risky industries — demand full disclosure or decline. Do not delegate the decision.}
	\end{frame}
	
	% ================================================================
	% SLIDE 5: TUTORIAL 4 — TRADE MISINVOICING
	% ================================================================
	\begin{frame}
		\frametitle{\fontsize{14}{17}\selectfont TUTORIAL 4: Trade Misinvoicing}
		
		\small
		\begin{block}{\fontsize{11}{13}\selectfont Topic: Trade-Based Money Laundering}
			\footnotesize
			Trade misinvoicing is a common money laundering technique. Understanding the indicators is essential for detection.
		\end{block}
		
		\vspace{0.03cm}
		
		\stepbox{\footnotesize
			\steptitle{Step 1:} Understand the scenario. A small textile exporter regularly sends shipments priced significantly below market rates. The foreign counterpart pays the remainder via separate wire transfers to a third party. \\
			\steptitle{Step 2:} Recognize what this suggests. This is \textbf{intentional trade mispricing for money laundering}. The exporter is undervaluing goods to move value out of the country. \\
			\steptitle{Step 3:} Understand how it works: \\
			\quad - Goods worth \$100,000 are invoiced at \$70,000 \\
			\quad - Buyer pays \$70,000 through normal channels \\
			\quad - Buyer sends \$30,000 through a separate wire to a third party \\
			\quad - Exporter has moved \$30,000 out of the country \\
			\steptitle{Step 4:} See the red flags. \textbf{Two indicators}: (1) Prices significantly below market rates, (2) Separate payments to third parties. Both suggest manipulation. \\
			\steptitle{Step 5:} Apply the lesson. When you see pricing that does not match market rates, look for other payments. The combination is a strong indicator of trade-based laundering.}
		
		\vspace{0.08cm}
		\tipbox{\footnotesize \textbf{CORE LESSON:} Trade misinvoicing is a common laundering technique. Look for prices that do not match market rates and separate payments to third parties.}
	\end{frame}
	
	% ================================================================
	% SLIDE 6: TUTORIAL 5 — INEFFECTIVE AML TRAINING
	% ================================================================
	\begin{frame}
		\frametitle{\fontsize{14}{17}\selectfont TUTORIAL 5: Ineffective AML Training}
		
		\small
		\begin{block}{\fontsize{11}{13}\selectfont Topic: AML Training Program Effectiveness}
			\footnotesize
			An AML training program is only effective if it produces measurable results. Certain indicators show the program is failing.
		\end{block}
		
		\vspace{0.03cm}
		
		\stepbox{\footnotesize
			\steptitle{Step 1:} Understand what indicates an ineffective training program. Two key indicators are: \\
			\quad (1) \textbf{Inconsistent application of CDD in client onboarding} \\
			\quad (2) \textbf{Employees demonstrate uncertainty when interviewed by auditors} \\
			\steptitle{Step 2:} Recognize why the first indicator matters. If employees apply Customer Due Diligence (CDD) \textbf{inconsistently}, they are not following procedures correctly. This creates gaps that criminals can exploit. \\
			\steptitle{Step 3:} Understand the second indicator. When employees are \textbf{uncertain when interviewed}, it shows they do not understand the policies. They may be guessing or making assumptions. \\
			\steptitle{Step 4:} See what is NOT the best answer. \textbf{Staff developing informal bypass workarounds} is a \textbf{consequence} of ineffective training, not a direct indicator. \\
			\steptitle{Step 5:} Apply the lesson. An effective training program produces \textbf{consistent, confident employees}. Inconsistency and uncertainty are red flags.}
		
		\vspace{0.08cm}
		\tipbox{\footnotesize \textbf{CORE LESSON:} Effective AML training produces consistent application of procedures and confident employees. Inconsistency and uncertainty indicate training failure.}
	\end{frame}
	
	% ================================================================
	% SLIDE 7: TUTORIAL 6 — SECOND LINE OF DEFENSE
	% ================================================================
	\begin{frame}
		\frametitle{\fontsize{14}{17}\selectfont TUTORIAL 6: Understanding the Lines of Defense}
		
		\small
		\begin{block}{\fontsize{11}{13}\selectfont Topic: AML Risk Management}
			\footnotesize
			Financial institutions have multiple lines of defense. Each line has a different role. Understanding who does what is essential.
		\end{block}
		
		\vspace{0.03cm}
		
		\stepbox{\footnotesize
			\steptitle{Step 1:} Understand the three lines of defense: \\
			\quad \textbf{First Line:} Operational management — they own and manage risks. They do the day-to-day work. \\
			\quad \textbf{Second Line:} Risk and compliance functions — they oversee and monitor risks. They ensure the first line is doing things correctly. \\
			\quad \textbf{Third Line:} Internal audit — they provide independent assurance. They check both lines. \\
			\steptitle{Step 2:} Understand what falls to the second line. The second line is responsible for \textbf{assessing whether new monitoring rules meet regulatory expectations}. \\
			\steptitle{Step 3:} See why this is the second line's role. The second line \textbf{oversees} the monitoring system. They do not design it (first line) and do not audit it (third line). They assess and challenge. \\
			\steptitle{Step 4:} Understand the boundary. The second line does NOT do operational work. They do NOT audit. They \textbf{oversee} and \textbf{challenge}. \\
			\steptitle{Step 5:} Apply the lesson. When determining responsibility, ask: Is this operational, oversight, or assurance?}
		
		\vspace{0.08cm}
		\tipbox{\footnotesize \textbf{CORE LESSON:} The second line of defense oversees and challenges. They assess whether monitoring rules meet regulatory expectations. They do not design or audit — they \textit{oversee}.}
	\end{frame}
	
	% ================================================================
	% SLIDE 8: TUTORIAL 7 — DATA INTEGRATION
	% ================================================================
	\begin{frame}
		\frametitle{\fontsize{14}{17}\selectfont TUTORIAL 7: Data Integration for Compliance}
		
		\small
		\begin{block}{\fontsize{11}{13}\selectfont Topic: Data Collection and Preparation}
			\footnotesize
			Integrating data from multiple sources into compliance platforms is a major challenge. Success requires standardization and alignment.
		\end{block}
		
		\vspace{0.03cm}
		
		\stepbox{\footnotesize
			\steptitle{Step 1:} Understand what ensures effective integration. Three key elements are: \\
			\quad (1) \textbf{Maintaining a common data dictionary across systems} \\
			\quad (2) \textbf{Mapping all incoming data to standardized field names} \\
			\quad (3) \textbf{Aligning control data requirements to system inputs} \\
			\steptitle{Step 2:} Recognize the importance of a common data dictionary. If different systems use different names for the same data, integration fails. A \textbf{common data dictionary} provides a shared language. \\
			\steptitle{Step 3:} Understand data mapping. \textbf{Mapping all incoming data to standardized field names} ensures that data from different sources can be compared and combined. \\
			\steptitle{Step 4:} See the role of alignment. \textbf{Aligning control data requirements to system inputs} ensures that the system receives the data it needs to perform compliance functions. \\
			\steptitle{Step 5:} Understand what is NOT correct. \textbf{Separating internal and external data into isolated systems} is the opposite of integration.}
		
		\vspace{0.08cm}
		\tipbox{\footnotesize \textbf{CORE LESSON:} Effective data integration requires a common dictionary, standardized mapping, and alignment of requirements. Separation of data is the enemy of integration.}
	\end{frame}
	
	% ================================================================
	% SLIDE 9: TUTORIAL 8 — CLUSTERING ANALYSIS
	% ================================================================
	\begin{frame}
		\frametitle{\fontsize{14}{17}\selectfont TUTORIAL 8: Clustering Analysis}
		
		\small
		\begin{block}{\fontsize{11}{13}\selectfont Topic: Data Collection — Clustering}
			\footnotesize
			Clustering analysis helps identify connections between accounts that might indicate criminal activity. Certain combinations strongly suggest meaningful relationships.
		\end{block}
		
		\vspace{0.03cm}
		
		\stepbox{\footnotesize
			\steptitle{Step 1:} Understand what clustering analysis does. It groups entities (accounts, customers, transactions) based on shared characteristics. The goal is to find relationships that might indicate criminal behavior. \\
			\steptitle{Step 2:} Recognize what suggests a meaningful connection. The combination of \textbf{same employer address}, \textbf{similar transaction values}, and \textbf{interlinked payments} strongly suggests a connection between accounts. \\
			\steptitle{Step 3:} Understand why this combination matters: \\
			\quad - Same employer address: Accounts are connected to the same organization \\
			\quad - Similar transaction values: They are doing similar business \\
			\quad - Interlinked payments: Money is flowing between them \\
			\steptitle{Step 4:} See the risk. This combination could indicate related parties, nominee accounts, or a structure designed to conceal ownership. \\
			\steptitle{Step 5:} Apply the lesson. When you see these three factors together, investigate further.}
		
		\vspace{0.08cm}
		\tipbox{\footnotesize \textbf{CORE LESSON:} Same employer address + similar transaction values + interlinked payments = meaningful connection. Investigate further.}
	\end{frame}
	
	% ================================================================
	% SLIDE 10: TUTORIAL 9 — PASSIVE LIVENESS DETECTION
	% ================================================================
	\begin{frame}
		\frametitle{\fontsize{14}{17}\selectfont TUTORIAL 9: Passive vs Active Liveness Detection}
		
		\small
		\begin{block}{\fontsize{11}{13}\selectfont Topic: Technology for Customer Onboarding}
			\footnotesize
			Liveness detection is used to verify that a person is physically present during identity verification. Passive and active methods differ significantly.
		\end{block}
		
		\vspace{0.03cm}
		
		\stepbox{\footnotesize
			\steptitle{Step 1:} Understand the distinction. \textbf{Active liveness} requires user interaction — asking the user to blink, turn head, smile, etc. \textbf{Passive liveness} requires minimal user interaction. \\
			\steptitle{Step 2:} Learn how passive liveness works. It \textbf{uses minimal user interaction and analyzes subtle facial movement or texture patterns}. The user does not need to perform specific actions. \\
			\steptitle{Step 3:} Understand the advantage of passive liveness. It is \textbf{smoother for the user} because they do not need to perform actions. It also detects sophisticated spoofing attempts by analyzing natural facial movements. \\
			\steptitle{Step 4:} Understand the active method. Active liveness requires the user to \textbf{perform specific actions} like blinking, smiling, or turning. This proves the person is present but can be less user-friendly. \\
			\steptitle{Step 5:} Apply the lesson. Passive liveness is more user-friendly because it requires minimal interaction. It analyzes natural facial patterns.}
		
		\vspace{0.08cm}
		\tipbox{\footnotesize \textbf{CORE LESSON:} Passive liveness detection uses minimal user interaction and analyzes subtle facial movements. It is more user-friendly than active methods.}
	\end{frame}
	
	% ================================================================
	% SLIDE 11: TUTORIAL 10 — ADDITIONAL ONBOARDING MEASURES
	% ================================================================
	\begin{frame}
		\frametitle{\fontsize{14}{17}\selectfont TUTORIAL 10: Additional Onboarding Measures}
		
		\small
		\begin{block}{\fontsize{11}{13}\selectfont Topic: Customer Onboarding Controls}
			\footnotesize
			Device recognition alone is not sufficient for customer authentication. Additional measures are needed to improve financial crime controls.
		\end{block}
		
		\vspace{0.03cm}
		
		\stepbox{\footnotesize
			\steptitle{Step 1:} Understand the scenario. A fintech startup relies solely on device recognition technology for customer authentication. The regulator is assessing this. \\
			\steptitle{Step 2:} Recognize the limitation. Device recognition identifies a device, not a person. A device can be shared or stolen. This is not sufficient for identity verification. \\
			\steptitle{Step 3:} Learn the additional measures needed. Two key measures are: \\
			\quad (1) \textbf{Implementing document verification with biometric authentication} \\
			\quad (2) \textbf{Integrating with national electronic identity verification systems} \\
			\steptitle{Step 4:} Understand why these are correct. \\
			\quad - Document verification + biometrics confirms the person is who they claim to be. \\
			\quad - National e-ID systems provide authoritative verification of identity. \\
			\steptitle{Step 5:} See what is NOT the best answer. \textbf{Enhancing digital footprint analysis through IP address verification} is helpful but not as reliable as document verification and national ID integration.}
		
		\vspace{0.08cm}
		\tipbox{\footnotesize \textbf{CORE LESSON:} Device recognition is not enough. Add document verification with biometrics and integrate with national identity systems.}
	\end{frame}
	
	% ================================================================
	% SLIDE 12: TUTORIAL 11 — CONCENTRATION RISK
	% ================================================================
	\begin{frame}
		\frametitle{\fontsize{14}{17}\selectfont TUTORIAL 11: Concentration Risk}
		
		\small
		\begin{block}{\fontsize{11}{13}\selectfont Topic: Money Laundering and Financial Crime}
			\footnotesize
			Concentration risk occurs when a financial institution has excessive exposure to a single counterparty, sector, or geographic area. Failing to address this risk has serious consequences.
		\end{block}
		
		\vspace{0.03cm}
		
		\stepbox{\footnotesize
			\steptitle{Step 1:} Understand what concentration risk is. It is the risk of loss from having too much exposure to a single counterparty, sector, or geographic area. \\
			\steptitle{Step 2:} Learn the consequences of failing to address concentration risk: \\
			\quad (1) \textbf{Increased exposure if a single counterparty defaults} \\
			\quad (2) \textbf{Decreased resilience in times of economic downturn} \\
			\steptitle{Step 3:} Understand why these are the correct consequences. If you have too much exposure to one counterparty, their default causes major losses. If the economy downturns, concentration amplifies the impact. \\
			\steptitle{Step 4:} See what is NOT correct. \textbf{Regulatory fines for poor credit modeling} is not a direct consequence of concentration risk. \textbf{Expansion into multiple business lines} reduces concentration risk. \\
			\steptitle{Step 5:} Apply the lesson. Diversification reduces concentration risk. Monitor exposures closely.}
		
		\vspace{0.08cm}
		\tipbox{\footnotesize \textbf{CORE LESSON:} Concentration risk creates vulnerability. If a single counterparty defaults or the economy downturns, the impact is amplified. Diversify to reduce risk.}
	\end{frame}
	
	% ================================================================
	% SLIDE 13: TUTORIAL 12 — BULK CASH AND DECENTRALIZED STORAGE
	% ================================================================
	\begin{frame}
		\frametitle{\fontsize{14}{17}\selectfont TUTORIAL 12: Bulk Cash and Decentralized Storage}
		
		\small
		\begin{block}{\fontsize{11}{13}\selectfont Topic: Money Laundering Vulnerabilities}
			\footnotesize
			Bulk cash movement and decentralized storage present significant vulnerabilities for money laundering. Understanding these risks is essential.
		\end{block}
		
		\vspace{0.03cm}
		
		\stepbox{\footnotesize
			\steptitle{Step 1:} Understand the vulnerability. \textbf{Vulnerabilities in bulk cash movement and decentralized storage methods} create opportunities for money laundering. \\
			\steptitle{Step 2:} Recognize why bulk cash is vulnerable. Cash can be moved across borders without detection. It does not leave an electronic trail. It can be stored in decentralized locations. \\
			\steptitle{Step 3:} Understand decentralized storage. Criminals can store cash in multiple locations: safe houses, warehouses, businesses. This makes it harder to detect and seize. \\
			\steptitle{Step 4:} See how criminals exploit this. They can move bulk cash across porous borders. They can store it in decentralized locations. They can then slowly introduce it into the financial system. \\
			\steptitle{Step 5:} Apply the lesson. Bulk cash movement and decentralized storage are significant vulnerabilities. Monitor for unexplained cash movements.}
		
		\vspace{0.08cm}
		\tipbox{\footnotesize \textbf{CORE LESSON:} Bulk cash can be moved across borders without detection and stored in decentralized locations. This creates significant money laundering vulnerabilities.}
	\end{frame}
	
	% ================================================================
	% SLIDE 14: TUTORIAL 13 — PERPETUAL KYC vs TRADITIONAL KYC
	% ================================================================
	\begin{frame}
		\frametitle{\fontsize{14}{17}\selectfont TUTORIAL 13: Perpetual KYC vs Traditional KYC}
		
		\small
		\begin{block}{\fontsize{11}{13}\selectfont Topic: Technology for Customer Onboarding}
			\footnotesize
			Perpetual KYC represents a fundamental shift from traditional periodic KYC. Understanding the difference is essential for modern compliance.
		\end{block}
		
		\vspace{0.03cm}
		
		\stepbox{\footnotesize
			\steptitle{Step 1:} Understand the distinction. \textbf{Perpetual KYC (PKYC)} triggers reviews based on real-time data changes. \textbf{Traditional KYC} reviews occur at fixed time intervals. \\
			\steptitle{Step 2:} See how PKYC works. PKYC monitors customer behavior continuously. When a triggering event occurs (new transaction pattern, change in ownership, etc.), the system updates the risk profile immediately. \\
			\steptitle{Step 3:} See how traditional KYC works. Traditional KYC reviews customers at set intervals: annually, biennially, etc. Changes between reviews are missed. \\
			\steptitle{Step 4:} Understand why PKYC is better. \\
			\quad - PKYC catches changes immediately \\
			\quad - PKYC reduces the window of exposure \\
			\quad - PKYC is more responsive to evolving risks \\
			\steptitle{Step 5:} Apply the lesson. PKYC is the future of customer due diligence. It is more effective than periodic reviews.}
		
		\vspace{0.08cm}
		\tipbox{\footnotesize \textbf{CORE LESSON:} Perpetual KYC triggers reviews based on real-time data changes. Traditional KYC reviews at fixed intervals. PKYC is more responsive and effective.}
	\end{frame}
	
	% ================================================================
	% SLIDE 14b: DATA GOVERNANCE COMMITTEE
	% ================================================================
	\begin{frame}
		\frametitle{\fontsize{14}{17}\selectfont TUTORIAL 14: Data Governance Committee}
		
		\small
		\begin{block}{\fontsize{11}{13}\selectfont Topic: Data Collection — Governance}
			\footnotesize
			A data governance committee plays a critical role in ensuring compliance data quality and trustworthiness.
		\end{block}
		
		\vspace{0.03cm}
		
		\stepbox{\footnotesize
			\steptitle{Step 1:} Understand the role of a data governance committee. It oversees the management of data across the organization. \\
			\steptitle{Step 2:} Learn what the committee does. \\
			\quad (1) \textbf{It strengthens trust in compliance data solutions through review processes}. \\
			\quad (2) \textbf{It supports compliance precision by reviewing key data outputs}. \\
			\steptitle{Step 3:} Understand why this matters. Compliance decisions are only as good as the data they are based on. The committee ensures the data is accurate and reliable. \\
			\steptitle{Step 4:} See the value of review processes. Regular review of data outputs catches errors before they affect compliance decisions. \\
			\steptitle{Step 5:} Apply the lesson. Data governance is essential for compliance. A committee provides oversight and accountability.}
		
		\vspace{0.08cm}
		\tipbox{\footnotesize \textbf{CORE LESSON:} A data governance committee strengthens trust in compliance data and supports precision through review processes.}
	\end{frame}
	
	% ================================================================
	% SLIDE 15: TUTORIAL 15 — DATA PREPARATION GOALS
	% ================================================================
	\begin{frame}
		\frametitle{\fontsize{14}{17}\selectfont TUTORIAL 15: Data Preparation Goals}
		
		\small
		\begin{block}{\fontsize{11}{13}\selectfont Topic: Data Collection — Preparation}
			\footnotesize
			Data preparation for compliance solutions has specific goals. Understanding these goals ensures successful implementation.
		\end{block}
		
		\vspace{0.03cm}
		
		\stepbox{\footnotesize
			\steptitle{Step 1:} Understand the goals of data preparation. Two key goals are: \\
			\quad (1) \textbf{Ensuring conformity to the data dictionary} \\
			\quad (2) \textbf{Verifying completeness and integrity of mapped fields} \\
			\steptitle{Step 2:} Recognize the importance of data dictionary conformity. If data does not conform to the dictionary, it cannot be integrated or compared. \\
			\steptitle{Step 3:} Understand completeness and integrity. \textbf{Completeness} means all required data is present. \textbf{Integrity} means the data is accurate and has not been altered. \\
			\steptitle{Step 4:} See what is NOT a goal. \textbf{Scheduling alerts every time compliance rules refresh} is not a data preparation goal — it is an operational activity. \\
			\steptitle{Step 5:} Apply the lesson. Data preparation focuses on making data ready for use. This means conforming to standards and ensuring quality.}
		
		\vspace{0.08cm}
		\tipbox{\footnotesize \textbf{CORE LESSON:} Data preparation ensures conformity to the data dictionary and verifies completeness and integrity of mapped fields.}
	\end{frame}
	
	% ================================================================
	% SLIDE 16: TUTORIAL 16 — ESCALATION AND REPORTING
	% ================================================================
	\begin{frame}
		\frametitle{\fontsize{14}{17}\selectfont TUTORIAL 16: Escalation and Reporting}
		
		\small
		\begin{block}{\fontsize{11}{13}\selectfont Topic: Data Collection — Escalation}
			\footnotesize
			When data matching identifies potential issues, escalation and reporting are essential next steps. Understanding the purpose is critical.
		\end{block}
		
		\vspace{0.03cm}
		
		\stepbox{\footnotesize
			\steptitle{Step 1:} Understand the purpose of escalation and reporting. The purpose is \textbf{to initiate enhanced due diligence and notify regulatory authorities if required}. \\
			\steptitle{Step 2:} Recognize why escalation matters. When data matching identifies a potential issue, action must be taken. Escalation ensures that the right people are notified. \\
			\steptitle{Step 3:} Understand enhanced due diligence. If an issue is identified, enhanced due diligence may be required. This means deeper investigation and review. \\
			\steptitle{Step 4:} Understand regulatory notification. If suspicious activity is identified, regulatory authorities may need to be notified through SARs. \\
			\steptitle{Step 5:} Apply the lesson. Escalation and reporting are not optional. They are required steps when issues are identified.}
		
		\vspace{0.08cm}
		\tipbox{\footnotesize \textbf{CORE LESSON:} Escalation and reporting initiate enhanced due diligence and notify regulatory authorities if required.}
	\end{frame}
	
	% ================================================================
	% SLIDE 17: TUTORIAL 17 — HASH FUNCTION IN BLOCKCHAIN
	% ================================================================
	\begin{frame}
		\frametitle{\fontsize{14}{17}\selectfont TUTORIAL 17: Hash Function in Blockchain}
		
		\small
		\begin{block}{\fontsize{11}{13}\selectfont Topic: Data Collection — Blockchain}
			\footnotesize
			Hash functions are a fundamental component of blockchain technology. Understanding what they do is essential for understanding blockchain-based compliance.
		\end{block}
		
		\vspace{0.03cm}
		
		\stepbox{\footnotesize
			\steptitle{Step 1:} Understand what a hash is. A hash is a mathematical function that converts input data of any size into a fixed-size string of characters. \\
			\steptitle{Step 2:} Understand the role of the hash in a blockchain block. The hash \textbf{acts as a unique block identifier}. \\
			\steptitle{Step 3:} Recognize why this is important. Each block in a blockchain has a unique hash. The hash is created from the block's data. If the data changes, the hash changes. \\
			\steptitle{Step 4:} Understand the security function. Because the hash is unique and changes if data changes, it provides \textbf{integrity}. It ensures that the block has not been tampered with. \\
			\steptitle{Step 5:} Apply the lesson. Hash functions provide integrity and uniqueness. They are essential to blockchain security.}
		
		\vspace{0.08cm}
		\tipbox{\footnotesize \textbf{CORE LESSON:} A hash in a blockchain block acts as a unique block identifier. It ensures data integrity by changing if the block data changes.}
	\end{frame}
	
	% ================================================================
	% SLIDE 18: SUMMARY
	% ================================================================
	\begin{frame}
		\frametitle{\fontsize{14}{17}\selectfont SUMMARY: Key Lessons from All Topics}
		
		\small
		\begin{block}{\fontsize{11}{13}\selectfont What We Have Learned:}
			\footnotesize
			\begin{enumerate}
				\item \textbf{DNFBP Risk Management:} Use client selection committees and standardized questionnaires.
				\item \textbf{Free-Trade Zones:} Lack of standardized inspection protocols creates money laundering risk.
				\item \textbf{TCSP Client Acceptance:} Decline clients with multiple red flags unless full disclosure provided.
				\item \textbf{Trade Misinvoicing:} Prices below market rates with separate payments = trade-based laundering.
				\item \textbf{AML Training:} Inconsistent CDD and employee uncertainty indicate training failure.
				\item \textbf{Second Line of Defense:} Assesses whether monitoring rules meet regulatory expectations.
				\item \textbf{Data Integration:} Common dictionary, standardized mapping, and alignment of requirements.
				\item \textbf{Clustering Analysis:} Same employer + similar values + interlinked payments = connection.
				\item \textbf{Liveness Detection:} Passive uses minimal interaction and analyzes subtle facial movements.
				\item \textbf{Onboarding Measures:} Add document verification with biometrics and national ID systems.
				\item \textbf{Concentration Risk:} Single counterparty default and economic downturn are key consequences.
				\item \textbf{Bulk Cash:} Movement and decentralized storage are significant vulnerabilities.
				\item \textbf{Perpetual KYC:} Triggers reviews based on real-time data changes.
				\item \textbf{Data Governance:} Committee strengthens trust and supports precision.
				\item \textbf{Data Preparation:} Conformity to dictionary and completeness/integrity of fields.
				\item \textbf{Escalation:} Initiates enhanced due diligence and notifies regulators if required.
				\item \textbf{Hash Function:} Acts as a unique block identifier in blockchain.
			\end{enumerate}
		\end{block}
		
		\vspace{0.05cm}
		\tipbox{\footnotesize \textbf{FINAL LESSON:} Every topic teaches us that understanding the \textit{underlying principle} is more important than memorizing facts. Know \textit{why} and you will know \textit{what} to do.}
	\end{frame}
	
	% ================================================================
	% SLIDE 19: THANK YOU
	% ================================================================
	\begin{frame}
		\centering
		\vspace{2.5cm}
		{\fontsize{22}{26}\selectfont \textbf{Thank You}} \\[0.3cm]
		{\fontsize{14}{17}\selectfont Keep learning. Keep growing.} \\[0.2cm]
		\rule{3cm}{1pt} \\[0.2cm]
		{\fontsize{10}{12}\selectfont \textit{Compliance is not about perfection — it is about continuous improvement.}}
		\vspace{0.5cm}
		{\fontsize{9}{11}\selectfont \textcolor{darkgray}{End of Tutorial}}
	\end{frame}
	
\end{document}